\documentclass[10pt,a4paper]{amsart}
\usepackage[margin=19mm]{geometry}
\usepackage{amsmath,amssymb,mathtools}
\usepackage[scr=rsfs,cal=euler]{mathalfa}
\usepackage{xfrac}
\usepackage[unicode,hidelinks,bookmarks=false]{hyperref}
\newcommand{\ie}{i.e.\ }
\newcommand{\cf}{cf.\ }
\newcommand{\resp}{respectively\ }
\newcommand{\Subsection}{\subsection}
\providecommand{\nobreakdash}{\nobreak\hbox{-}}
\setlength{\emergencystretch}{2em}
\multlinegap0pt
\usepackage{bm}
\usepackage{bbm}
\usepackage{dsfont}
\usepackage{xspace}
\usepackage{cancel}
\usepackage{mathtools}
\usepackage{amsthm}
\makeatletter
\@ifundefined{thm}{\newtheorem{thm}{Thm}}{}
\@ifundefined{lemma}{\newtheorem{lemma}[thm]{Lemma}}{}
\makeatother
\newcommand{\eps}{\varepsilon}
\title[A simple method of computing Riemann-type integrals]{A simple method of computing Riemann-type integrals}
\author{Luciano Tubaro}
\date{Source: arXiv:2609.05562v2 (CC BY 4.0)}
\begin{document}
\maketitle

\noindent\emph{Source and licence.} A simple method of computing Riemann-type integrals. Version arXiv:2609.05562v2 (CC BY 4.0), distributed under the Creative Commons Attribution 4.0 International licence, as recorded in the arXiv licence metadata for this exact version and shown on the abstract page https://arxiv.org/abs/2609.05562v2. Full rights notice: licenses/arxiv-2609.05562-CC-BY-4.0.txt.\par\smallskip

\noindent\textit{Editorial scope.} Counted unit: the Lemma whose source label is \texttt{lem:l2} in the article (source0.tex lines 1765--1789, source label \texttt{lem:l2}), delivered together with the article's own complete original proof of it (source0.tex lines 1791--1930, one \texttt{proof} environment of 140 lines). No mathematics is written, repaired or re-derived: statement and proof are the article's own text line for line. Required context retained uncounted: lem:l1 (lines 1525--1532). Every definition, display and label of those lines is retained unchanged. Omitted from this package and not redistributed: 1--1524, 1533--1764, 1790--1790, 1931--2462. Nothing inside a retained excerpt is omitted or altered: every retained body is delivered exactly as the article's lines are written, so no omission receipt of any kind accompanies this package. Citations: the retained excerpts carry no citation command at all, so no bibliography entry of the article is transcribed and no reference is redistributed. Reference adaptation: none was needed and none was made; every reference inside the delivered unit resolves inside it, so no unresolved reference or citation is produced. Printed slips kept literal: no printed word, symbol, hypothesis or number of the counted statement or of its proof was altered. Layout and class: the article's own class is \texttt{article} with packages including amsmath, amsfonts, amssymb, calrsfs, amsthm; the wrapper is the lane's pinned \texttt{amsart} editorial wrapper at 10pt on A4, which restates the theorem-like environments the retained text needs and adds the article's own macro definitions that the retained text needs (\texttt{\textbackslash eps}), with extra wrapper packages (bm, bbm, dsfont, xspace, cancel, mathtools). The delivered numbering is the unit's own. Assets: no retained span contains a figure environment, a graphics-inclusion command, a PDF-page inclusion, a table or a TikZ picture; no image file is copied into this package, the package declares no asset dependency and no image file is redistributed with it. Licence: version arXiv:2609.05562v2 of this article is distributed under the Creative Commons Attribution 4.0 International licence, as recorded in the arXiv licence metadata for this exact version and shown on the abstract page https://arxiv.org/abs/2609.05562v2. A full rights notice is delivered at licenses/arxiv-2609.05562-CC-BY-4.0.txt. Characters: every retained excerpt is pure ASCII, so the delivered engine, XeLaTeX with its default Latin Modern text font, reports no missing character.
\begin{lemma}\label{lem:l1}
The sums $\sum_{i=0}^{n-1}W_{t_i}^k(W_{t_{i+1}}-W_{t_i})^\ell$ converge, for
$\ell\ge3$, to $0$ in $L^2(\Omega)$; for $\ell=2$,
\[
\sum_{i=0}^{n-1}W_{t_i}^k(W_{t_{i+1}}-W_{t_i})^2\longrightarrow\int_a^bW_t^k\,dt
\qquad\text{in }L^2(\Omega).
\]
\end{lemma}

\begin{lemma}\label{lem:l2}
Let $X_t$, $t\in[a,b]$, have continuous paths (a.s.), and let $A_t$,
$t\in[a,b]$, be a continuous, non-decreasing process such that, for
every $s\le t$ in $[a,b]$ and every partition $\pi$ \emph{containing}
$s$ and $t$ as partition points (i.e.\ $s,t\in\{t_0,\ldots,t_n\}$, so
that the sub-sum below is unambiguous, with no boundary fragment left
over at either end),
\[
\sum_{i\,:\,s\le t_i<t_{i+1}\le t}(X_{t_{i+1}}-X_{t_i})^2
\longrightarrow A_t-A_s\qquad\text{in probability, as }\delta(\pi)\to0
\]
along such $s,t$-containing partitions (taking $s=a$, $t=b$ gives, in
particular, $\sum_{i=0}^{n-1}(X_{t_{i+1}}-X_{t_i})^2\to A_b-A_a$, and
every partition trivially contains $a,b$ as its endpoints). Then, for
$\ell\ge3$,
\[
\sum_{i=0}^{n-1}X_{t_i}^k(X_{t_{i+1}}-X_{t_i})^\ell\longrightarrow0
\qquad\text{in probability},
\]
while for $\ell=2$,
\[
\sum_{i=0}^{n-1}X_{t_i}^k(X_{t_{i+1}}-X_{t_i})^2\longrightarrow\int_a^bX_t^k\,dA_t
\qquad\text{in probability}.
\]
\end{lemma}

\begin{proof}
Throughout, write $\Delta X_i:=X_{t_{i+1}}-X_{t_i}$ for the increments of a
generic partition $\pi\colon a=t_0<\cdots<t_n=b$. We do \emph{not} use a
semimartingale decomposition $X=Y+M$, nor any martingale property of $A$:
the conclusion follows directly from the two stated hypotheses ---
convergence of $\sum(\Delta X_i)^2$ to $A_t-A_s$ on every sub-interval
$[s,t]$, along partitions containing $s,t$ as grid points, and continuity
of the path $t\mapsto X_t$ --- together with the
monotonicity of $A$, used below. This is the same elementary
Riemann--Stieltjes convergence argument used, in the deterministic
setting, throughout this note, now carried out ``in probability'' instead
of pointwise, and it requires nothing from the general theory of
semimartingales.

\medskip\noindent\textit{Step 1: the case $\ell=2$ for a simple weight.}
Let $\phi$ be a (non-random, or more generally $\mathcal F_a$-measurable)
step function on $[a,b]$, $\phi(t)=c_j$ for $t\in[s_j,s_{j+1})$, for a
fixed partition $a=s_0<s_1<\cdots<s_N=b$, and set
$M_\phi:=\max_j|c_j|<\infty$. For $\pi$ \emph{containing} $\{s_j\}$
(i.e.\ $\{s_j\}\subset\{t_0,\ldots,t_n\}$),
\begin{align*}
\sum_i\phi(t_i)(\Delta X_i)^2&=\sum_{j=0}^{N-1}c_j\!\!
\sum_{i\,:\,t_i\in[s_j,s_{j+1})}\!\!(\Delta X_i)^2\\
&\ \xrightarrow[\delta(\pi)\to0]{P}\ \sum_{j=0}^{N-1}c_j(A_{s_{j+1}}-A_{s_j})
=:\int_a^b\phi(t)\,dA_t,
\end{align*}
using the hypothesis on each of the finitely many fixed sub-intervals
$[s_j,s_{j+1}]$ and the fact that a finite sum of terms, each converging
in probability, converges in probability. This proves the $\ell=2$
statement for a step-function weight in place of $X_{t_i}^k$, for
partitions containing $\{s_j\}$, directly from the hypothesis.

\smallskip\noindent\emph{Extension to arbitrary $\pi$.} The hypothesis
of the Lemma, and hence the display above, is stated only for partitions
containing the fixed points $\{s_j\}$; we now remove this restriction,
for this same fixed $\phi$, since it is needed for what follows. Given
an arbitrary partition $\pi\colon a=t_0<\cdots<t_n=b$ (not necessarily
containing $\{s_j\}$), let $\widehat\pi:=\pi\cup\{s_j\}$ be the partition
obtained by inserting the (at most $N-1$ interior) points of $\{s_j\}$
not already in $\pi$; then $\widehat\pi\supset\{s_j\}$, and inserting
points can only shrink cells, so $\delta(\widehat\pi)\le\delta(\pi)$.
At most $N-1$ cells of $\pi$ are split by this insertion (one per
interior point $s_j$); every other cell of $\pi$ is a cell of
$\widehat\pi$ too, contributing identically to both sums below.
For a split cell $[t_i,t_{i+1}]\ni s_j$ (in its interior), the
$\pi$-sum contributes $\phi(t_i)(X_{t_{i+1}}-X_{t_i})^2$, while the
$\widehat\pi$-sum contributes
$\phi(t_i)(X_{s_j}-X_{t_i})^2+\phi(s_j)(X_{t_{i+1}}-X_{s_j})^2$; both
are, in modulus, at most $M_\phi\,\delta(\pi)^2$ (each factor
$|X_{t_{i+1}}-X_{t_i}|,|X_{s_j}-X_{t_i}|,|X_{t_{i+1}}-X_{s_j}|$ is at
most $\delta(\pi)$, being a sub-increment of a cell of $\pi$ of length
$\le\delta(\pi)$), so their difference is at most $3M_\phi\,\delta(\pi)^2$.
Summing over the at most $N-1$ affected cells,
\[
\Big|\sum_i\phi(t_i)(\Delta X_i)^2-\sum_i\phi(\hat t_i)(\Delta\widehat
X_i)^2\Big|\le3(N-1)M_\phi\,\delta(\pi)^2\longrightarrow0
\]
\emph{deterministically}, as $\delta(\pi)\to0$ -- with $N$ fixed (it
depends only on $\phi$, chosen before $\pi$), this bound has nothing
random or probabilistic left in it, unlike every other estimate in this
proof. Since $\widehat\pi\supset\{s_j\}$ and
$\delta(\widehat\pi)\le\delta(\pi)\to0$, the first display of this Step
applies to $\widehat\pi$, giving
$\sum_i\phi(\hat t_i)(\Delta\widehat X_i)^2\to\int_a^b\phi\,dA_t$ in
probability as $\delta(\pi)\to0$; combined with the deterministic bound
above,
\[
\sum_i\phi(t_i)(\Delta X_i)^2\ \xrightarrow[\delta(\pi)\to0]{P}\
\int_a^b\phi(t)\,dA_t
\]
for \emph{every} partition $\pi$, not only those containing $\{s_j\}$.
This is the form of Step~1 used in Step~2 below.

\medskip\noindent\textit{Step 2: the case $\ell=2$ for $X_t^k$, by
approximation.} Since $X$ has continuous paths a.s.\ and $[a,b]$ is
compact, the (random) modulus of continuity
\[
\omega_{X^k}(\eta):=\sup_{\substack{s,t\in[a,b]\\|t-s|\le\eta}}|X_t^k-X_s^k|
\]
satisfies $\omega_{X^k}(\eta)\to0$ a.s.\ as $\eta\to0^+$ (pathwise
uniform continuity on the compact $[a,b]$), hence also in probability.
So, given $\eps>0$ and $\rho>0$, there is a \emph{deterministic}
$\eta=\eta(\eps,\rho)>0$ --- not depending on $\omega$ --- such that
\[
P\big(\omega_{X^k}(\eta)>\eps\big)<\rho.
\]
Fix, once and for all, a deterministic partition $a=s_0<\cdots<s_N=b$ of
mesh $<\eta$ (independent of $\omega$), and define the step function
$\phi(t):=X_{s_j}^k$ for $t\in[s_j,s_{j+1})$; $\phi$ is random through
its values $X_{s_j}^k$, but its underlying partition $\{s_j\}$ is fixed,
so Step~1 (in the form just established for arbitrary partitions)
applies to it directly, without restricting $\pi$. On the event
$\{\omega_{X^k}(\eta)\le\eps\}$, of probability $\ge1-\rho$,
$\sup_{t\in[a,b]}|X_t^k-\phi(t)|\le\eps$ pathwise. For \emph{every}
partition $\pi$, on this event,
\begin{align*}
\Big|\sum_iX_{t_i}^k(\Delta X_i)^2-\sum_i\phi(t_i)(\Delta X_i)^2\Big|
&\le\Big(\sup_t|X_t^k-\phi(t)|\Big)\sum_i(\Delta X_i)^2\\
&\le\eps\sum_i(\Delta X_i)^2,
\end{align*}
and likewise $\big|\int_a^bX_t^k\,dA_t-\int_a^b\phi(t)\,dA_t\big|
\le\eps(A_b-A_a)$, since $A$ is non-decreasing. Since $\sum_i(\Delta
X_i)^2\to A_b-A_a$ in probability (the hypothesis, $t=b$), it is bounded
in probability: given $\rho>0$ there is $M$ with
$P\big(\sum_i(\Delta X_i)^2>M\big)<\rho$ for all $\pi$ with $\delta(\pi)$
small enough. On the intersection of this event with
$\{\omega_{X^k}(\eta)\le\eps\}$ --- of probability $\ge1-2\rho$, for the
deterministic $\eta=\eta(\eps,\rho)$ fixed above --- both error terms are
at most $\eps M$ and $\eps(A_b-A_a)$ respectively; combined with Step~1
applied to this fixed, deterministic $\phi$, a standard ``$3\eps$''
argument gives
\[
\sum_iX_{t_i}^k(\Delta X_i)^2\ \xrightarrow[\delta(\pi)\to0]{P}\ \int_a^bX_t^k\,dA_t,
\]
which is the $\ell=2$ statement, established directly from the hypothesis
and the continuity of $X$, without any martingale argument. Working
throughout with a fixed, deterministic partition $\{s_j\}$ --- rather
than a partition constructed after fixing $\omega$ --- keeps the
approximation argument and the ``in probability'' statement cleanly
separated.

\medskip\noindent\textit{Step 3: the case $\ell\ge3$.} By Step~2 (with
$|X_t|^k$ in place of $X_t^k$), $\sum_i|X_{t_i}|^k(\Delta X_i)^2\to
\int_a^b|X_t|^k\,dA_t$ in probability, so this quantity is bounded in
probability (tight) as $\delta(\pi)\to0$. Since $X$ is (a.s.) uniformly
continuous on $[a,b]$, $\eta(\pi):=\max_i|\Delta X_i|\to0$ a.s.; then, for
$\ell\ge3$,
\[
\Big|\sum_iX_{t_i}^k(\Delta X_i)^\ell\Big|
\le\eta(\pi)^{\ell-2}\sum_i|X_{t_i}|^k(\Delta X_i)^2,
\]
a product of a term $\to0$ a.s.\ (hence in probability) and a tight term,
hence itself $\to0$ in probability. (As in Lemma~\ref{lem:l1}, this
argument gives convergence in probability for the net $\delta(\pi)\to0$;
an almost-sure statement is \emph{not} established by this route and
would need either a fixed sequence of partitions with
$\sum_n\delta(\pi_n)<\infty$ together with a Borel--Cantelli argument, or
the $L^2$-theory for the quadratic variation of a bounded martingale
together with passage to a subsequence.)
\end{proof}

\end{document}
